% filas del capítulo 07: símbolo & significado & unidad \\ (4.ª columna, macro, opcional; toda fila termina en \\)
$N_{\mathrm U}$, $N_{\mathrm M}$ & nocional en \USD{} y en \MXN{} del cross currency swap, $N_{\mathrm M}=N_{\mathrm U}X_{\mathrm p}$ & \USD; \MXN \\
$X_{\mathrm p}$ & tipo de cambio pactado (\texttt{fx\_rate}) & \MXN{} por \USD \\
$\psi$, $\psi^{*}$ & diferencial sobre la pata TIIE-F y su valor par & decimal anual (1~\pb$=10^{-4}$) \\
$\theta$ & signo del lado: $+1$ si el tenedor paga la pata en dólares (\texttt{PayUsd}), $-1$ si la paga en pesos & adimensional \\
$e_0$, $e_i$ & inicio del periodo 1 y día de pago $i$, en días desde la fecha de valuación & días \\
$\tau_i$ & devengo ACT/360 del periodo $i$, $(e_i-e_{i-1})/360$ & años \\
$P_c(e)$ & factor de descuento de la curva $c$ a $e$ días ($1$ si $e\le0$) & adimensional \\
$G_{\mathrm U}$, $G_{\mathrm M}$ & factores acumulados de las fijaciones conocidas del periodo 1 de cada pata & adimensional \\
$\mathbf 1_{\mathrm{ini}}$, $\mathbf 1_{\mathrm{fin}}$ & indicadores de intercambio de nocional inicial y final & adimensional \\
$V_{\mathrm U}$, $V_{\mathrm M}$ & valor de cada pata visto por quien la recibe & \USD; \MXN \\
$\Ann_{\mathrm M}$ & anualidad de la pata en pesos, $\sum_i\tau_iP_{\cMXNcol}(e_i)$ & años \\
$k_i$, $\ell$, $i^{\circ}$ & días de pago desde el inicio, días transcurridos y primer periodo pendiente & días; días; índice \\
